\documentclass[a4paper]{arrowhead}

\usepackage[yyyymmdd]{datetime}
\usepackage{etoolbox}
\usepackage[T1]{fontenc}
\usepackage[utf8]{inputenc}
\usepackage{multirow}
\usepackage{longtable}
\usepackage{amssymb}
\usepackage{float}
\usepackage{tikz}
\usetikzlibrary{positioning,arrows.meta,fit,calc,shapes.misc,shapes.geometric}

\renewcommand{\dateseparator}{-}
\setlength{\emergencystretch}{3em}
\lstset{basicstyle=\ttfamily\footnotesize,breaklines=true,columns=fullflexible,frame=single,rulecolor=\color{gray!50},backgroundcolor=\color{gray!6}}
\tikzset{
  ahsys/.style={draw=ArrowheadBlue,very thick,fill=ArrowheadBlue!6,rounded corners=2pt,minimum width=3.2cm,minimum height=1.1cm,align=center,font=\small\bfseries},
  ahext/.style={draw=ArrowheadGrey,thick,dashed,fill=gray!5,rounded corners=2pt,minimum width=3.2cm,minimum height=1.0cm,align=center,font=\small},
  ahsrv/.style={draw=ArrowheadPurple,thick,fill=ArrowheadPurple!8,rounded corners=6pt,minimum height=0.55cm,inner xsep=4pt,align=center,font=\scriptsize\ttfamily},
  ahcsm/.style={draw=ArrowheadTeal,thick,fill=ArrowheadTeal!8,rounded corners=6pt,minimum height=0.55cm,inner xsep=4pt,align=center,font=\scriptsize\ttfamily},
  ahdb/.style={draw=ArrowheadOrange,thick,fill=ArrowheadOrange!8,cylinder,shape border rotate=90,aspect=0.25,minimum height=0.9cm,align=center,font=\scriptsize},
  ahflow/.style={-{Stealth[length=2.2mm]},thick,ArrowheadGrey},
  ahlabel/.style={font=\scriptsize,fill=white,inner sep=1pt,align=center},
}

\newcommand{\fparam}[1]{\textit{\textcolor{ArrowheadBlue}{#1}}}
%% Special references
\newcommand{\mref}[1]{{\textcolor{ArrowheadPurple}{\hyperref[sec:model:#1]{#1}}}}
\newcommand{\oref}[1]{{\textcolor{ArrowheadBlue}{\hyperref[sec:operations:#1]{#1}}}}
\newcommand{\pdef}[1]{{\textcolor{ArrowheadGrey}{#1\label{sec:model:primitives:#1}\label{sec:model:primitives:#1s}}}}
\newcommand{\pref}[1]{{\textcolor{ArrowheadGrey}{\hyperref[sec:model:primitives:#1]{#1}}}}

\newrobustcmd\osubsection[5]{
  \addtocounter{subsection}{1}
  \addcontentsline{toc}{subsection}{\protect\numberline{\thesubsection}operation \textcolor{ArrowheadBlue}{#1}}
  \subsection*{
    \thesubsection\quad
    #2 \textcolor{ArrowheadPurple}{#3} \\
    \small
    \hspace*{0.075\textwidth}\begin{minipage}{0.1\textwidth}
      \vspace*{1mm}
      Operation: \\
      \notblank{#4}{Input: \\}{}
      \notblank{#5}{Output: \\}{}
    \end{minipage}
    \begin{minipage}{0.825\textwidth}
      \vspace*{1mm}
      \textcolor{ArrowheadBlue}{#1} \\
      \notblank{#4}{\mref{#4} \\}{}
      \notblank{#5}{\mref{#5} \\}{}
    \end{minipage}
  }
  \label{sec:operations:#1}
}
\newrobustcmd\msubsection[2]{
  \addtocounter{subsection}{1}
  \addcontentsline{toc}{subsection}{\protect\numberline{\thesubsection}#1 \textcolor{ArrowheadPurple}{#2}}
  \subsection*{\thesubsection\quad#1 \textcolor{ArrowheadPurple}{#2}}
  \label{sec:model:#2} \label{sec:model:#2s}
}

\begin{document}

%% Arrowhead Document Properties
\ArrowheadTitle{workflowManagement HTTP/TLS/JSON}
\ArrowheadServiceID{workflowManagement}
\ArrowheadType{Interface Design Description}
\ArrowheadTypeShort{IDD}
\ArrowheadVersion{5.4.0}
\ArrowheadDate{2026-11-23}
\ArrowheadAuthor{Documentation maintenance}
\ArrowheadStatus{REVIEW DRAFT}
\ArrowheadContact{Repository maintainers}
\ArrowheadFooter{\href{www.arrowhead.eu}{www.arrowhead.eu}}
\ArrowheadSetup
%%

%% Front Page
\begin{center}
  \vspace*{1cm}
  \huge{\arrowtitle}

  \vspace*{0.2cm}
  \LARGE{\arrowtype}
  \vspace*{1cm}

  \Large{Service ID: \textit{``\arrowid''}}
  \vspace*{\fill}

  

  \vspace*{1cm}
  \vspace*{\fill}

  % Front Page Abstract
  \begin{abstract}
    This document describes the HTTP/TLS/JSON interface of the \texttt{workflowManagement} service of the Eclipse Arrowhead workflow support systems, service version 1.0.0, documentation edition 5.4.0. It realises the abstract operations of the \texttt{workflowManagement} SD.
  \end{abstract}

  \vspace*{1cm}
\end{center}

\newpage
%%

%% Table of Contents
\tableofcontents
\newpage
%%

\section{Overview}
\label{sec:overview}
This document describes how the \texttt{workflowManagement} service, which gives operators the worklog, the known executors and the work items, is realised with HTTP~\cite{rfc9110}
over TLS~\cite{rfc8446} and JSON~\cite{rfc8259} payloads. Request and response bodies use content type
\texttt{application/json}; JSON field names are case-sensitive and fields not listed in this document must not be sent
(the recipe decoder rejects unknown fields). Optional fields may be omitted; absent and \texttt{null} are equivalent.

The service is offered by the WorkflowManager on its Arrowhead connector (default port 8482) with base path \texttt{/\allowbreak{}workflowmanager/\allowbreak{}mgmt}.
The ServiceRegistry entry lists every operation with its HTTP method and path relative to the base path.

The rest of this document is organized as follows. Section~\ref{sec:operations} specifies each operation,
Section~\ref{sec:model} the JSON data models, Section~\ref{sec:security} the security requirements,
Section~\ref{sec:errors} the error representation and Section~\ref{sec:implementations} the implementations.

\newpage
\section{Service Operations}
\label{sec:operations}

{\small
\begin{longtable}{|p{3.0cm}|p{1.5cm}|p{5.3cm}|p{3.0cm}|p{3.0cm}|}

\hline
\rowcolor{gray!33} Operation & Method & Path & Input & Output \\ \hline
\endhead
\oref{query-worklog} & GET & \texttt{/\allowbreak{}workflowmanager/\allowbreak{}mgmt/\allowbreak{}worklog} & \mref{WorklogFilter} & \mref{WorklogList} \\ \hline
\oref{query-executors} & GET & \texttt{/\allowbreak{}workflowmanager/\allowbreak{}mgmt/\allowbreak{}executors} & -- & \mref{ExecutorList} \\ \hline
\oref{query-work-items} & GET & \texttt{/\allowbreak{}workflowmanager/\allowbreak{}mgmt/\allowbreak{}work-\allowbreak{}items} & \mref{WorklogFilter} & \mref{WorkItemList} \\ \hline
\end{longtable}
}

\osubsection{query-worklog}{GET}{/workflowmanager/mgmt/worklog}{WorklogFilter}{WorklogList}
Returns worklog entries, optionally of one session.

{\small
\begin{longtable}{|p{3.0cm}|p{1.4cm}|p{3.2cm}|p{8.2cm}|}

\hline
\rowcolor{gray!33} Parameter & In & Type & Description \\ \hline
\endhead
\fparam{sessionId} & query & \pref{SessionId} & Optional session. \\ \hline
\end{longtable}
}


{\small
\begin{longtable}{|p{2.0cm}|p{13.8cm}|}

\hline
\rowcolor{gray!33} Status & When \\ \hline
\endhead
Success & 200 OK. \\ \hline
401 & No valid identity token, or the token is bound to another certificate (\texttt{AUTH}). \\ \hline
403 & The consumer has no ConsumerAuthorization grant for this operation (\texttt{FORBIDDEN}). \\ \hline
\end{longtable}
}

\osubsection{query-executors}{GET}{/workflowmanager/mgmt/executors}{}{ExecutorList}
Returns the executor set $WE$ with $cap(w)$: the providers of \texttt{workflowExecution} in the ServiceRegistry plus statically configured executors.

{\small
\begin{longtable}{|p{2.0cm}|p{13.8cm}|}

\hline
\rowcolor{gray!33} Status & When \\ \hline
\endhead
Success & 200 OK. \\ \hline
401 & No valid identity token, or the token is bound to another certificate (\texttt{AUTH}). \\ \hline
403 & The consumer has no ConsumerAuthorization grant for this operation (\texttt{FORBIDDEN}). \\ \hline
\end{longtable}
}

\osubsection{query-work-items}{GET}{/workflowmanager/mgmt/work-items}{WorklogFilter}{WorkItemList}
Returns work items, optionally of one session.

{\small
\begin{longtable}{|p{3.0cm}|p{1.4cm}|p{3.2cm}|p{8.2cm}|}

\hline
\rowcolor{gray!33} Parameter & In & Type & Description \\ \hline
\endhead
\fparam{sessionId} & query & \pref{SessionId} & Optional session. \\ \hline
\end{longtable}
}


{\small
\begin{longtable}{|p{2.0cm}|p{13.8cm}|}

\hline
\rowcolor{gray!33} Status & When \\ \hline
\endhead
Success & 200 OK. \\ \hline
401 & No valid identity token, or the token is bound to another certificate (\texttt{AUTH}). \\ \hline
403 & The consumer has no ConsumerAuthorization grant for this operation (\texttt{FORBIDDEN}). \\ \hline
\end{longtable}
}


\newpage
\section{Data Models}
\label{sec:model}
The JSON objects exchanged by the operations, in alphabetic order. Field types refer to other data models or to the
primitives of Section~\ref{sec:model:primitives}.

\msubsection{enum}{DispatchMode}
How a step reaches an executor.
 Encoded as a JSON string.
{\small
\begin{longtable}{|p{3.9cm}|p{11.9cm}|}

\hline
\rowcolor{gray!33} Value & Meaning \\ \hline
\endhead
\texttt{push} & The WorkflowManager selects an executor and calls \texttt{workflowExecution.start} (default). \\ \hline
\texttt{pull} & The WorkflowManager publishes the work item; an executor claims it with \texttt{workflowWorkItem.claim}. \\ \hline
\texttt{event} & Reserved for EventHandler notification; currently handled as \texttt{pull}. \\ \hline
\end{longtable}
}

\msubsection{struct}{ErrorMessage}
Error response of every operation (common library \texttt{ErrorMessageDTO}).

{\small
\begin{longtable}{|p{3.7cm}|p{4.2cm}|p{1.6cm}|p{6.3cm}|}

\hline
\rowcolor{gray!33} Field & Type & Mandatory & Description \\ \hline
\endhead
\texttt{errorMessage} & \pref{String} & yes & Explanation. \\ \hline
\texttt{errorCode} & \pref{Integer} & yes & HTTP status code. \\ \hline
\texttt{exceptionType} & \pref{String} & yes & \texttt{INVALID\_PARAMETER}, \texttt{AUTH}, \texttt{FORBIDDEN}, \texttt{DATA\_NOT\_FOUND}, \texttt{EXTERNAL\_SERVER\_ERROR}, \ldots \\ \hline
\texttt{origin} & \pref{String} & no & The operation that failed, e.g. \texttt{POST /workflowmanager/work-item/report}. \\ \hline
\end{longtable}
}

\msubsection{struct}{ExecutorEndpoint}
An executor $w \in WE$ with $cap(w)$.

{\small
\begin{longtable}{|p{3.7cm}|p{4.2cm}|p{1.6cm}|p{6.3cm}|}

\hline
\rowcolor{gray!33} Field & Type & Mandatory & Description \\ \hline
\endhead
\texttt{name} & \pref{SystemName} & yes & Executor. \\ \hline
\texttt{scheme} & \pref{String} & yes & \texttt{https} or \texttt{http}. \\ \hline
\texttt{address} & \pref{String} & yes & Address. \\ \hline
\texttt{port} & \pref{Integer} & yes & Port. \\ \hline
\texttt{basePath} & \pref{String} & yes & Base path of \texttt{workflowExecution}. \\ \hline
\texttt{capabilities} & \pref{List}$<$\pref{String}$>$ & yes & Capabilities. \\ \hline
\texttt{services} & \pref{List}$<$\pref{ServiceDefinition}$>$ & yes & Runnable services. \\ \hline
\texttt{workstation} & \pref{String} & no & Workstation. \\ \hline
\texttt{mode} & \mref{DispatchMode} & yes & Mode. \\ \hline
\texttt{source} & \pref{String} & yes & \texttt{registry} or \texttt{static}. \\ \hline
\end{longtable}
}

\msubsection{struct}{ExecutorList}
Executors known to the WorkflowManager.

{\small
\begin{longtable}{|p{3.7cm}|p{4.2cm}|p{1.6cm}|p{6.3cm}|}

\hline
\rowcolor{gray!33} Field & Type & Mandatory & Description \\ \hline
\endhead
\texttt{entries} & \pref{List}$<$\mref{ExecutorEndpoint}$>$ & yes & Executors. \\ \hline
\texttt{count} & \pref{Integer} & yes & Number of entries. \\ \hline
\end{longtable}
}

\msubsection{struct}{WorkItemList}
Work items.

{\small
\begin{longtable}{|p{3.7cm}|p{4.2cm}|p{1.6cm}|p{6.3cm}|}

\hline
\rowcolor{gray!33} Field & Type & Mandatory & Description \\ \hline
\endhead
\texttt{entries} & \pref{List}$<$\mref{WorkItemRecord}$>$ & yes & Work items. \\ \hline
\texttt{count} & \pref{Integer} & yes & Number of entries. \\ \hline
\end{longtable}
}

\msubsection{struct}{WorkItemRecord}
Book-keeping of one work item.

{\small
\begin{longtable}{|p{3.7cm}|p{4.2cm}|p{1.6cm}|p{6.3cm}|}

\hline
\rowcolor{gray!33} Field & Type & Mandatory & Description \\ \hline
\endhead
\texttt{workItemId} & \pref{WorkItemId} & yes & Identifier. \\ \hline
\texttt{sessionId} & \pref{SessionId} & yes & Session. \\ \hline
\texttt{action} & \pref{Name} & yes & Action. \\ \hline
\texttt{step} & \pref{Name} & yes & Step. \\ \hline
\texttt{executor} & \pref{SystemName} & no & Assigned or claiming executor. \\ \hline
\texttt{attempt} & \pref{Integer} & yes & Attempt. \\ \hline
\texttt{state} & \mref{WorkItemState} & yes & State. \\ \hline
\texttt{message} & \pref{String} & no & Last message. \\ \hline
\texttt{createdAt} & \pref{DateTime} & yes & Created. \\ \hline
\texttt{updatedAt} & \pref{DateTime} & yes & Updated. \\ \hline
\end{longtable}
}

\msubsection{enum}{WorkItemState}
State of a work item at the WorkflowManager.
 Encoded as a JSON string.
{\small
\begin{longtable}{|p{3.9cm}|p{11.9cm}|}

\hline
\rowcolor{gray!33} Value & Meaning \\ \hline
\endhead
\texttt{OPEN} & Published for claiming (pull). \\ \hline
\texttt{DISPATCHED} & Sent to an executor (push). \\ \hline
\texttt{CLAIMED} & Claimed by an executor (pull). \\ \hline
\texttt{BOUND} & The executor reported a provider. \\ \hline
\texttt{DONE} & Completed. \\ \hline
\texttt{FAILED} & Failed. \\ \hline
\texttt{ABORTED} & Aborted. \\ \hline
\end{longtable}
}

\msubsection{struct}{WorklogEntry}
One audit entry of the WorkflowManager.

{\small
\begin{longtable}{|p{3.7cm}|p{4.2cm}|p{1.6cm}|p{6.3cm}|}

\hline
\rowcolor{gray!33} Field & Type & Mandatory & Description \\ \hline
\endhead
\texttt{id} & \pref{Integer} & yes & Identifier. \\ \hline
\texttt{createdAt} & \pref{DateTime} & yes & Time. \\ \hline
\texttt{sessionId} & \pref{SessionId} & no & Session, if any. \\ \hline
\texttt{recipe} & \pref{String} & no & Recipe, if any. \\ \hline
\texttt{step} & \pref{Name} & no & Step, if any. \\ \hline
\texttt{message} & \pref{String} & yes & Event. \\ \hline
\texttt{detail} & \pref{String} & no & Detail. \\ \hline
\end{longtable}
}

\msubsection{struct}{WorklogFilter}
Optional filter.

{\small
\begin{longtable}{|p{3.7cm}|p{4.2cm}|p{1.6cm}|p{6.3cm}|}

\hline
\rowcolor{gray!33} Field & Type & Mandatory & Description \\ \hline
\endhead
\texttt{sessionId} & \pref{SessionId} & no & Only entries/items of this session. \\ \hline
\end{longtable}
}

\msubsection{struct}{WorklogList}
Worklog entries.

{\small
\begin{longtable}{|p{3.7cm}|p{4.2cm}|p{1.6cm}|p{6.3cm}|}

\hline
\rowcolor{gray!33} Field & Type & Mandatory & Description \\ \hline
\endhead
\texttt{entries} & \pref{List}$<$\mref{WorklogEntry}$>$ & yes & Entries. \\ \hline
\texttt{count} & \pref{Integer} & yes & Number of entries. \\ \hline
\end{longtable}
}


\subsection{Primitives}
\label{sec:model:primitives}
The abstract primitive types of the SD are represented in JSON~\cite{rfc8259} as follows.
\subsubsection*{alias \pdef{DateTime}}
JSON string in the RFC 3339 format~\cite{rfc3339}, e.g. \texttt{2026-11-23T10:15:30Z}.

\subsubsection*{alias \pdef{Integer}}
JSON number without fraction.

\subsubsection*{alias \pdef{List}}
JSON array.

\subsubsection*{alias \pdef{Name}}
JSON string matching \texttt{\^{}[A-Za-z][A-Za-z0-9\_-]\{0,62\}\$}.

\subsubsection*{alias \pdef{ServiceDefinition}}
JSON string.

\subsubsection*{alias \pdef{SessionId}}
JSON number (64-bit integer); a path segment when used in a path.

\subsubsection*{alias \pdef{String}}
JSON string (UTF-8).

\subsubsection*{alias \pdef{SystemName}}
JSON string.

\subsubsection*{alias \pdef{WorkItemId}}
JSON string holding a UUID; a path segment when used in a path.



\section{Security}
\label{sec:security}
The interface is offered only over TLS~\cite{rfc8446} (TLS 1.3 or 1.2) with mutual authentication: the consumer must
present an Arrowhead certificate issued under the local cloud's certificate authority, and certificates on the
cloud's CRL are refused. The service is registered with interface template \texttt{generic\_https} and interface
policy \texttt{CERT\_AUTH}.

Every request must carry the consumer's identity token, obtained from the Authentication system and bound to the
consumer's certificate:
\begin{lstlisting}
Authorization: Bearer IDENTITY-TOKEN//<identity token>
\end{lstlisting}
A token issued to another system's certificate is refused (401). The provider then checks a ConsumerAuthorization
grant for the operation (\texttt{authorization.verify}, \texttt{service.authorization.enabled=true}); a missing grant is
refused with 403. Sysop only (management policy); a ConsumerAuthorization grant is not sufficient.


\section{Errors}
\label{sec:errors}
Every error response carries an \mref{ErrorMessage} body with content type \texttt{application/json}, for example:
\begin{lstlisting}
{ "errorMessage": "Work item 5d2c0b1e-... is held by another executor",
  "errorCode": 403, "exceptionType": "FORBIDDEN",
  "origin": "POST /workflowmanager/work-item/report" }
\end{lstlisting}
Status codes follow HTTP semantics~\cite{rfc9110}: 400 invalid parameter, 401 authentication, 403 forbidden, 404 not
found, 409 conflict with the current state, 500 internal error and 503 failure of an external system.


\section{Implementations}
\label{sec:implementations}
\begin{itemize}
\item \textbf{Java (reference).} \texttt{Eclipse-Arrowhead-java/ah5-Workflow-java}, module \texttt{workflow-manager}. Paths, methods and service
operations are defined once in \texttt{workflow-model} (\texttt{wire/WorkflowServices.java}, \texttt{wire/ServiceOperations.java}) and checked by the provider-side authorization filter.
\item \textbf{Go.} Not implemented: \texttt{Eclipse-Arrowhead-go} has no workflow systems. A Go implementation must
provide exactly the paths, methods, JSON field names and status codes of this IDD to be interoperable.
\item \textbf{Rust.} Not implemented: \texttt{Eclipse-Arrowhead-rust} has no workflow systems. The same requirement applies.
\end{itemize}
Until the Go and Rust implementations exist, the workflow systems are \emph{not} interoperable across the three
implementations; only the Java systems interoperate with Java, Go and Rust core systems through the common Arrowhead
interfaces (ServiceRegistry, Authentication, ConsumerAuthorization, DynamicServiceOrchestration).


\newpage
\bibliographystyle{IEEEtran}
\bibliography{bibliography}

\newpage

\section{Revision History}
\subsection{Amendments}

\noindent\begin{tabularx}{\textwidth}{| p{1cm} | p{2.2cm} | p{1.6cm} | X | p{3.6cm} |} \hline
\rowcolor{gray!33} No. & Date & Version & Subject of Amendments & Author \\ \hline
1 & 2026-11-23 & \arrowversion & Initial edition: interface design description of workflowManagement over HTTP/TLS/JSON, written from the Java reference implementation (ah5-Workflow-java). & Documentation maintenance \\ \hline
\end{tabularx}

\subsection{Quality Assurance}

\noindent\begin{tabularx}{\textwidth}{| p{1cm} | p{2.2cm} | p{1.6cm} | X |} \hline
\rowcolor{gray!33} No. & Date & Version & Approved by \\ \hline
1 & & \arrowversion & Pending review \\ \hline
\end{tabularx}

\end{document}
